\chapter{Perspectivas \'eticas, limitaciones, conclusiones y trabajo futuro}
\label{cap:conclusiones}

\section{Perspectivas \'eticas, sociales y culturales}
El desarrollo de artefactos tecnol\'ogicos orientados a comunidades hist\'oricamente vulneradas exige una reflexi\'on \'etica que trascienda la optimizaci\'on de indicadores f\'isicos. El primer compromiso es ecol\'ogico: el motor incorpora la pausa con $\text{HR} < 30.0\,\%$ y el tope de extracci\'on del $33.0\,\%$ de la niebla disponible para que la instalaci\'on no compita por humedad con la flora xerof\'itica ni altere los microclimas nocturnos. Ese compromiso tiene un costo que debe hacerse expl\'icito, ya que el tope explica la mayor parte del agua que las reglas dejan de captar frente al temporizador (\cref{sec:comparacion_modos}), y por eso su valor es una decisi\'on que corresponde discutir con las comunidades y no solo con criterios t\'ecnicos.

El segundo compromiso es la gobernanza territorial. Cabo de la Vela pertenece al resguardo ind\'igena de la Alta y Media Guajira, y ning\'un piloto de campo puede concebirse sin un proceso aut\'entico de consulta previa, libre e informada, conforme al Convenio 169 de la Organizaci\'on Internacional del Trabajo y a la jurisprudencia constitucional, incluida la Sentencia T-302 de 2017 \parencite{corte2017t302}. La tecnolog\'ia debe subordinarse a las estructuras tradicionales de decisi\'on de las autoridades wayuu. A ello se suma el riesgo de extractivismo epist\'emico: el informe y el tablero est\'an redactados en castellano, con resumen en ingl\'es, y la ausencia de una versi\'on en \emph{wayuunaiki} es una limitaci\'on \'etica evidente, de modo que cualquier fase experimental deber\'a incorporar la traducci\'on y validaci\'on de los conceptos y alertas del colector con hablantes de la comunidad. Por \'ultimo, el agua cosechada debe considerarse un bien com\'un esencial para la vida, la salud y la alimentaci\'on en la rancher\'ia, lo que implica que el mantenimiento sea autogestivo y no genere nuevas dependencias comerciales externas.

\section{Limitaciones del estudio}
La limitaci\'on de fondo es el nivel de madurez: el artefacto es un modelo anal\'itico y una arquitectura de control evaluada en simulaci\'on, en el rango TRL 2--3, sin validaci\'on en banco de pruebas ni prototipos de campo. De ella se derivan las dem\'as. Las cin\'eticas de remoci\'on de sal, fatiga y degradaci\'on de la membrana no proceden de literatura experimental de captadores, sino de supuestos de orden de magnitud (\cref{cap:marco_teorico}), y el uso de la impedancia como sustituto de la sal acumulada no se ha validado en mallas o membranas de niebla (\cref{sec:antecedentes}). Tampoco se ha simulado la falla de sensores, de la que depende la garant\'ia de seguridad.

A estas limitaciones del modelo se suman las amenazas a la validez de la comparaci\'on discutidas en el \cref{sec:validez}. El temporizador modelado carece de corte t\'ermico y de pausa ecol\'ogica, por lo que una parte de la ventaja en desgaste corresponde a las salvaguardas y no a la l\'ogica reactiva; la recuperaci\'on se mide en momentos distintos en cada modo, lo que limita la comparabilidad del veredicto sobre recuperaci\'on; y la banda de transici\'on se deriva y se verifica con la misma semilla. Las corridas, adem\'as, cubren 24 horas con una sola noche, ocho nodos y un clima sint\'etico con ciclos arm\'onicos y ruido gaussiano, sin series meteorol\'ogicas multianuales ni fen\'omenos como El Ni\~no o La Ni\~na.

Por \'ultimo, el impacto comunitario es una estimaci\'on simplificada. La proyecci\'on mensual supone un mes de 30 d\'ias id\'enticos, sin p\'erdidas de almacenamiento ni de distribuci\'on y sin estacionalidad de la niebla, y la cifra del ICBF es una entrega institucional por comunidad, no la dotaci\'on total de agua de una rancher\'ia. En el plano operativo, el tablero depende de la disponibilidad de la API para mostrar las series temporales.

\section{Conclusiones}
En atenci\'on a los objetivos propuestos y a cada parte de la pregunta de investigaci\'on, las conclusiones se formulan dentro de los l\'imites del modelo y de sus supuestos.

\textbf{a) En cuanto a la reducci\'on del desgaste}, el control basado en reglas se degrada menos que el temporizador en todas las condiciones evaluadas, con reducciones de la degradaci\'on acumulada entre el $\SI{61.5}{\percent}$ y el $\SI{100}{\percent}$, y la ventaja se mantiene al menos hasta una fatiga de $2.0\,\%$ por pulso. Sin embargo, una parte importante de esa ventaja proviene de que el temporizador modelado no tiene corte t\'ermico; aun descontando su p\'erdida t\'ermica, la reducci\'on se mantiene entre el $52.7\,\%$ y el $80.0\,\%$ en los escenarios analizados, lo que sugiere que la l\'ogica basada en condici\'on aporta una ventaja propia.

\textbf{b) En cuanto a la recuperaci\'on hidr\'aulica}, la ventaja del control reactivo es condicional y est\'a regida por la incrustaci\'on diaria: aparece dentro de la banda de transici\'on definida en el \cref{sec:banda}, con cruces interpolados entre $1633$ y $1867\;\Omega\text{/d\'ia}$, y se mantiene por encima de ella. Con los flujos de niebla de proyectos reales, el control reactivo solo supera al temporizador en el $\SI{4.8}{\percent}$ de los casos, 3 de 62 muestras de incertidumbre, de modo que en la mayor\'ia de los sitios su beneficio ser\'ia la protecci\'on de la membrana y no un aumento del agua recuperada.

\textbf{c) En cuanto a la seguridad operativa y los l\'imites ecol\'ogicos}, la prioridad de las reglas R1 y R2 asegura por construcci\'on, dentro del modelo y suponiendo sensores sin falla, que ninguna decisi\'on supere $45\;\si{\celsius}$ ni extraiga agua con $\text{HR} < 30\,\%$. En las simulaciones esto se tradujo en cero p\'erdida t\'ermica, cero violaciones ecol\'ogicas frente a ocho del temporizador y menor consumo de energ\'ia en todos los escenarios.

\textbf{d) En cuanto al impacto comunitario y el dimensionamiento}, abastecer a una comunidad de 50 personas con el est\'andar b\'asico de la OMS ($20\;\text{L/persona/d\'ia}$) demanda entre 100 y 332 nodos de $1\;\text{m}^2$ con los flujos calibrados en la corrida de referencia. Frente a la entrega mensual de $1000\;\text{L}$ por comunidad que reporta el ICBF, una bater\'ia de 3 a 13 nodos podr\'ia igualarla, lo que la sit\'ua como un posible complemento local para el hogar, sin sustituir la obligaci\'on estatal de garantizar infraestructura permanente de agua potable.

\noindent\textbf{Conclusi\'on general.}
Dentro del modelo y con sus supuestos, el motor de reglas determinista de Axioma-H ofrece una arquitectura de control segura por construcci\'on, auditable regla por regla y con menor consumo energ\'etico que el temporizador modelado. Su beneficio principal ser\'ia prolongar la vida \'util de la membrana al evitar ciclos t\'ermicos innecesarios y excursiones por encima de $45\;\si{\celsius}$, mientras que la ganancia en agua recuperada depende del sitio. Antes de afirmar estas ventajas fuera de la simulaci\'on, es necesario medir la eficacia y la fatiga del pulso y comparar con un temporizador dotado de las mismas salvaguardas.

\section{Trabajo futuro (Ciclo 3)}
\label{sec:trabajo_futuro}
La continuidad de la investigaci\'on se organiza en torno a dos hip\'otesis derivadas y contrastables, adem\'as de la hip\'otesis H1-Ciclo 3 (\cref{sec:hipotesis}). La primera, \textbf{H2}, sostiene que, frente a un temporizador dotado del mismo corte a $45\;\si{\celsius}$ y de la misma pausa ecol\'ogica, el control reactivo mantiene una menor degradaci\'on acumulada. Su contraste exige una corrida que hoy el artefacto no ofrece: un modo de temporizador con salvaguardas, simulado sobre el mismo entorno y con los mismos niveles de salinidad del barrido. La hip\'otesis quedar\'ia refutada si en esa corrida el temporizador con salvaguardas alcanzara una degradaci\'on igual o menor que la de las reglas en alg\'un nivel de salinidad, lo que indicar\'ia que la ventaja en desgaste proviene de las salvaguardas y no del control basado en condici\'on.

La segunda, \textbf{H3}, sostiene que la banda de transici\'on de incrustaci\'on diaria se conserva en una validaci\'on fuera de muestra, es decir, con semillas y cin\'eticas de pulso no usadas para derivarla. Su contraste consiste en repetir la robustez de la respuesta con semillas distintas de la 42 y de las semillas 1 a 3, y con eficacias del pulso fuera de la cuadr\'icula original, y comparar los nuevos cruces con la banda. La hip\'otesis quedar\'ia refutada si, en esas corridas, aparecieran puntos con ventaja positiva claramente por debajo de la banda o negativa claramente por encima, o si los cruces interpolados cayeran fuera de ella.

Estas hip\'otesis se complementan con mediciones de laboratorio y de campo. En laboratorio, la prioridad es medir sobre muestras de membrana la fracci\'on de sal que remueve un pulso de $90\;\text{s}$ a $44.0\;\si{\celsius}$, de la que dependen las conjeturas 1 y 2, y la fatiga por pulso mediante ciclos repetidos, que debe contrastarse con el criterio de refutaci\'on de H1. En Cabo de la Vela, una estaci\'on de monitoreo deber\'ia registrar el rendimiento de niebla con un colector est\'andar y la incrustaci\'on diaria, para ubicar el sitio respecto de la banda de transici\'on. Queda, adem\'as, una pregunta abierta previa a todas ellas: ?`es la impedancia superficial un sensor v\'alido de la incrustaci\'on en una malla o membrana de niebla, y con qu\'e electrodos o superficie conductora puede medirse sin alterar la captaci\'on? La respuesta condiciona tanto el dise\~no del motor como la interpretaci\'on de la carga salina del modelo.

En el plano tecnol\'ogico y social, el trabajo futuro incluye acoplar un m\'odulo neuronal ligero para predecir frentes de niebla a corto plazo, manteniendo el motor de reglas como guardarra\'il de decisi\'on; estimar el costo nivelado del litro cosechado frente al acarreo por carrotanque; y desarrollar con docentes y l\'ideres wayuu una versi\'on biling\"ue de la interfaz de supervisi\'on, como parte del proceso de consulta previa.
